\pdfminorversion=7
\documentclass[a4paper,final]{article}
\usepackage[14pt]{extsizes}
\usepackage[T1]{fontenc}
\usepackage[utf8]{inputenc}
\usepackage[english]{babel}

\addto\captionsenglish{\renewcommand{\figurename}{Fig.}}

% --- Math ---
\usepackage{amsmath,amssymb,amsfonts}
\usepackage{nccmath}

% --- Page geometry ---
\usepackage[left=20mm, top=20mm, right=20mm, bottom=20mm, footskip=10mm]{geometry}

% --- Text/layout helpers ---
\usepackage{ragged2e}
\usepackage{setspace}
\usepackage{indentfirst}

% --- Graphics ---
\usepackage{graphicx}
\DeclareGraphicsExtensions{.jpg,.png,.pdf}
\usepackage{float}
\usepackage{wrapfig}
\usepackage{tikz}
\usepackage{pgfplots}
\pgfplotsset{compat=1.18}

% --- Captions ---
\usepackage[font=small,singlelinecheck=false,justification=centering,labelsep=period]{caption}

% --- Lists ---
\usepackage{enumitem}

% --- Hyperlinks ---
\usepackage{hyperref}
\hypersetup{colorlinks=true, linkcolor=blue, urlcolor=blue}

% --- Code listings ---
\usepackage{xcolor}
\usepackage{listings}
\usepackage{pdfpages}
\definecolor{codegreen}{rgb}{0,0.6,0}
\definecolor{codegray}{rgb}{0.5,0.5,0.5}
\definecolor{codepurple}{rgb}{0.58,0,0.82}
\definecolor{backcolour}{rgb}{0.95,0.95,0.92}

\lstdefinestyle{mystyle}{
	backgroundcolor=\color{backcolour},
	commentstyle=\color{codegreen},
	keywordstyle=\color{magenta},
	numberstyle=\tiny\color{codegray},
	stringstyle=\color{codepurple},
	basicstyle=\ttfamily\footnotesize,
	breaklines=true,
	breakatwhitespace=false,
	captionpos=b,
	keepspaces=true,
	numbers=left,
	numbersep=5pt,
	showspaces=false,
	showstringspaces=false,
	showtabs=false,
	tabsize=4,
	columns=fullflexible,
	extendedchars=true,
	inputencoding=utf8,
	literate=
	{•}{{-}}1
	{─}{{-}}1
	{–}{{-}}1
	{—}{{-}}1
	{→}{{$\to$}}1
	{≠}{{$\neq$}}1
	{³}{{$^3$}}1
	{ü}{{\"u}}1
}

\lstset{style=mystyle}
\renewcommand{\lstlistingname}{Listing}
\newcommand{\coderef}[2]{\hyperref[#2]{\texttt{\detokenize{#1}}}}
\DeclareUnicodeCharacter{202F}{\,}
\sloppy

% --- Other document tweaks ---
\renewcommand{\arraystretch}{1.4}
\parindent=24pt
\parskip=0pt
\tolerance=2000
\flushbottom

\addto\captionsenglish{\def\refname{References}}

\begin{document}
	\begin{center}
		\normalsize\bf{{MINISTRY OF SCIENCE AND HIGHER EDUCATION OF THE RUSSIAN FEDERATION}

			\hfill \break
			Federal State Autonomous Educational Institution of Higher Education\\ "Peter the Great St. Petersburg Polytechnic University"\\[10pt]}
		\hfill \break
		\normalsize{Institute of Computer Science and Cybersecurity}\\[10pt]

		\normalsize{Higher School of Technology and Artificial Intelligence}\\[10pt]

		\normalsize{Program: 02.03.01 Mathematics and Computer Science}\\
		\vspace{50pt}

		\vspace{1pt}
		\begin{center}
Coursework report\\
on the discipline "Graph Theory"\\
Implementation of graph theory algorithms\\
for laboratory works No. 1--5

		\end{center}

		\hfill \break
		\hfill \break
	\end{center}
	\small{ \bf
		\begin{tabular}{lrrl}
			\!\!\!Student, & \hspace{0cm} & & \\
			\!\!\!group 5130201/40001 & \hspace{1.5cm} & \underline{\hspace{3cm}} & \hspace{-0.3cm}Ovi Md Shamin Yasir  \\\\
			\!\!\!Teacher, & \hspace{1cm} & \\
			\!\!\! & \hspace{0.1cm} & \underline{\hspace{3cm}} &  Vostrov Alexey Vladimirovich\\\\
			&&\hspace{3.5cm}
		\end{tabular}
		\begin{flushright}
			"\underline{\hspace{1cm}}"\underline{\hspace{2.5cm}} 2026
		\end{flushright}
	}

	\hfill \break
	\hfill \break
	\begin{center} \small\bf{St. Petersburg, 2026} \end{center}
	\thispagestyle{empty}
	% END OF TITLE PAGE

\newpage
\renewcommand{\contentsname}{Contents}
\tableofcontents

\newpage
\section*{Introduction}
\addcontentsline{toc}{section}{Introduction}

\indent This report describes 5 laboratory works. Each work includes the implementation of different algorithms on randomly generated connected acyclic graphs. The graphs are generated using the \textit{normal distribution} according to Vadzinsky's handbook~\cite{vadzinsky}.

\medskip
\noindent\textbf{LABORATORY WORK No. 1}

\noindent 1.\; Randomly generate a connected acyclic directed graph according to the \textit{normal distribution} with a user-defined number of vertices. The distribution parameters are constants selected experimentally. The distribution is taken from Vadzinsky's handbook~\cite{vadzinsky}. Generation is performed by vertex degrees.

\noindent 2.\; Calculate vertex eccentricities, determine the center of the graph and the diametral vertices.

\noindent 3.\; Implement Shimbell's method for the weight matrix generated by the specified distribution. The user enters the number of edges. The answer must contain the minimum and/or maximum path represented as a matrix. The program must support generating only positive, only negative, and mixed weights at the user's choice.

\noindent 4.\; Determine whether a route can be built from one specified vertex to another, where both vertices are entered by the user, and output the number of such routes.

\noindent\textit{Variant: normal distribution.}

\medskip
\noindent\textbf{LABORATORY WORK No. 2}

\noindent 1.\; For the given graphs randomly generated in the previous work, perform \textit{depth-first graph traversal}.

\noindent 2.\; Generate a weight matrix, either positive or containing negative weights at the user's choice, and find the shortest path between two user-selected vertices using the \textit{Floyd-Warshall algorithm}. The result includes the distance matrix and the path itself as a sequence of vertices.

\noindent 3.\; Compare the running speed of the implemented algorithms by the number of iterations.

\noindent\textit{Variant: (c) depth-first graph traversal; (j) Floyd-Warshall algorithm.}

\medskip
\noindent\textbf{LABORATORY WORK No. 3}

\noindent 1.\; Based on the generated directed graph, randomly obtain the capacity and cost matrices.

\noindent 2.\; For the obtained graph, find the maximum flow using the \textit{Ford-Fulkerson algorithm}.

\noindent 3.\; Compute the specified minimum-cost flow, where the required flow value is $\lfloor 2/3 \cdot \mathrm{max} \rfloor$ and $\mathrm{max}$ is the maximum flow. Use the previously implemented Floyd-Warshall algorithm.

\medskip
\noindent\textbf{LABORATORY WORK No. 4}

\noindent 1.\; For the given undirected graphs randomly generated in the first work, find the number of spanning trees using Kirchhoff's matrix-tree theorem.

\noindent 2.\; Build a minimum-weight spanning tree for the generated undirected weighted graph using \textit{Boruvka's algorithm}. Encode the obtained spanning tree with the Prufer code and decode it while preserving the edge weights.

\noindent 3.\; Implement \textit{minimum edge cover} construction on the original graph and on the obtained spanning tree at the user's choice.

\noindent\textit{Variant: (b) Boruvka's algorithm; (g) minimum edge cover.}

\medskip
\noindent\textbf{LABORATORY WORK No. 5}

\noindent 1.\; For the given undirected graphs randomly generated in the first work, check whether the graph is Eulerian. If it is not, modify the graph and log what was changed. Build an Eulerian cycle.

\noindent 2.\; Build the shortest spanning tree using Boruvka's algorithm from the fourth work. Based on this spanning tree, obtain the \textit{fundamental cutset system} and print it. Implement obtaining all other cutsets from the fundamental ones using the symmetric difference operation; the selected cutsets are entered by the user.

\noindent\textit{Variant: 36, even -- fundamental cutset system.}

\medskip
The laboratory works were implemented in C++ in the CLion development environment.

\newpage
	\section{Mathematical Description}
	\subsection{Graph}

A simple graph $G(V, E)$ is a combination of two sets: a non-empty set $V$ and a set $E$ of unordered pairs of different elements of $V$. The set $V$ is called the set of vertices, and the set $E$ is called the set of edges.
\[
G(V,E) = \langle V, E\rangle,\quad V \neq \emptyset,\quad E \subseteq V \times V,\quad \{v,v\} \notin E,\quad v \in V,
\]
that is, the set $E$ consists of 2-element subsets of the set $V$.

Related terms:
\begin{itemize}
  \item A vertex of a graph $G$ is an element of the vertex set $v \in V(G)$;
  \item An edge of a graph $G$ is an element of the edge set $e \in E(G)$, or $e = \{v_1, v_2\}$, where $v_1 \in V(G)$ and $v_2 \in V(G)$;
  \item The elements of a graph $G$ are its vertices $v \in V(G)$ and edges $e \in E(G)$;
  \item A vertex $v$ is incident to an edge $e$ if $v \in e$; in this case, $e$ is also said to be an edge at $v$;
  \item Adjacent vertices are vertices $v_1$ and $v_2$ such that $\{v_1,v_2\} \in E(G)$;
  \item The degree of a vertex $d(v)$ is the number of edges incident to it.
\end{itemize}

	\subsection{Normal Distribution}

The normal, or Gaussian, distribution is a continuous probability distribution. Its probability density is
\[
f(x) = \frac{1}{\sigma\sqrt{2\pi}}\,\exp\!\left(-\frac{(x-\mu)^2}{2\sigma^2}\right),
\]
where $\mu$ is the expected value and $\sigma^2$ is the variance.

To generate random vertex degrees, an approximate method from Vadzinsky's handbook~\cite{vadzinsky} is used. A random variable $X$ with normal distribution $\mathcal{N}(\mu,\sigma^2)$ is approximated as follows:
\[
X \approx \sqrt{\frac{12}{n}}\left(\sum_{i=1}^{n} r_i - \frac{n}{2}\right),
\]
where $r_i \sim U(0,1)$ are independent random numbers uniformly distributed on $[0,1]$, and $n$ is the number of summands, which is the approximation accuracy parameter. The distribution is taken from the handbook on probability distributions~\cite{vadzinsky}.

	\subsection{Eccentricity, Center and Diametral Vertices}

The \textbf{eccentricity} of a vertex $v$ in a graph $G$ is the maximum distance from $v$ to the other vertices:
\[
\varepsilon(v) = \max_{u \in V}\, d(v, u),
\]
where $d(v, u)$ is the length of the shortest path between vertices $v$ and $u$.

The \textbf{radius} of a graph is the minimum eccentricity among all vertices:
\[
\mathrm{rad}(G) = \min_{v \in V}\, \varepsilon(v).
\]

The \textbf{diameter} of a graph is the maximum eccentricity:
\[
\mathrm{diam}(G) = \max_{v \in V}\, \varepsilon(v).
\]

The \textbf{center} of a graph is the set of vertices with minimum eccentricity equal to the radius:
\[
Z(G) = \{v \in V \mid \varepsilon(v) = \mathrm{rad}(G)\}.
\]

\textbf{Diametral} vertices are vertices whose eccentricity is equal to the diameter of the graph. Eccentricities are computed using breadth-first search (BFS) from each vertex in time $O(V(V+E))$.

	\subsection{Shimbell's Method}

Shimbell's method makes it possible to find minimum or maximum paths between vertices that consist of a specified number of edges.

\begin{itemize}
  \item The operation of multiplying two values $a$ and $b$ while raising a matrix to a power corresponds to their algebraic sum, that is:
  \[a * b = b * a \;\to\; a + b = b + a,\qquad a * 0 = 0 * a = 0 \;\to\; a + 0 = 0.\]
  \item The operation of adding two values $a$ and $b$ is replaced by selecting the maximum or minimum of these values, that is:
  \[a + b = b + a = \min(\max)\,\{a,\,b\}.\]
  Zeros are ignored in this operation.
  \[w_{ij} = \min\bigl(\max(w_{i1}+w_{1j}),\ldots,(w_{in}+w_{nj})\bigr).\]
\end{itemize}

	\subsection{Depth-First Search Traversal}

Depth-first search is one of the basic graph traversal methods. It is often used to check connectivity, find cycles and strongly connected components, and perform topological sorting.

The general idea of the algorithm is as follows: for each unvisited vertex, all unvisited adjacent vertices must be found, and the search must be repeated for them.

\begin{enumerate}
  \item Choose any vertex among the unvisited vertices and denote it by $u$;
  \item Run the procedure $\mathrm{dfs}(u)$:
  \begin{itemize}
    \item Mark vertex $u$ as visited;
    \item For each unvisited vertex adjacent to $u$, denoted by $v$, run $\mathrm{dfs}(v)$;
  \end{itemize}
  \item Repeat steps 1 and 2 until all vertices are visited.
\end{enumerate}

The running time of the algorithm is estimated as $O(V+E)$.

	\subsection{Floyd-Warshall Algorithm}

The Floyd-Warshall algorithm is intended for finding shortest paths between all pairs of vertices in a weighted graph, including graphs with negative edge weights, provided there are no negative cycles.

Let $d^{(k)}[i][j]$ denote the length of the shortest path from vertex $i$ to vertex $j$ that passes only through intermediate vertices from the set $\{1, 2, \ldots, k\}$. The recurrence relation is
\[
d^{(k)}[i][j] = \min\!\left(d^{(k-1)}[i][j],\; d^{(k-1)}[i][k] + d^{(k-1)}[k][j]\right).
\]

The initial conditions are
\[
d^{(0)}[i][j] = \begin{cases} w(i,j), & \text{if edge } (i,j) \text{ exists},\\ \infty, & \text{otherwise}. \end{cases}
\]

As a result, $d^{(n)}[i][j]$ contains the length of the shortest path from $i$ to $j$. The answer contains the distance matrix and the path itself as a sequence of vertices. The complexity of the algorithm is $O(V^3)$.

	\subsection{Ford-Fulkerson Algorithm}

The Ford-Fulkerson algorithm solves the problem of finding the maximum flow in a network. The idea of the algorithm is as follows:

\begin{itemize}
  \item Initially, the flow value is set to 0 for all $u$ and $v$ in the transport network.
  \item Then the flow is iteratively increased by searching for an augmenting path from the source $s$ to the sink $t$.
  \item The process is repeated while an augmenting path can be found.
\end{itemize}

For the Edmonds-Karp implementation used in the program, where augmenting paths are found by BFS, the running time is $O(VE^2)$.

\medskip
\noindent\textbf{The algorithm for finding a specified minimum-cost flow} can use one of the shortest-path algorithms from the source to the sink. In this work, the required total flow is routed by repeatedly finding the shortest path by cost in the residual network. The objective function is
\[
\min\!\left(\sum_{(u,v)\in E} c(u,v)\cdot f(u,v)\right),
\]
where $E$ is the set of all edges in the graph, $c(u,v)$ is the cost of flow on edge $(u,v)$, and $f(u,v)$ is the flow on edge $(u,v)$. In the implemented version, the shortest path in the residual cost graph is found using Floyd-Warshall.

	\subsection{Spanning Tree}

Let $G = (V, E)$ be a graph. A spanning subgraph of a graph $G = (V, E)$ is a subgraph containing all vertices. A spanning subgraph that is a tree is called a spanning tree. A disconnected graph has no spanning tree. A connected graph may have many spanning trees.

A minimum spanning tree of a graph $G$ is a spanning tree $T = (V, E_T)$ of the graph $G$ whose sum of edge weights is minimal.

	\subsection{Kirchhoff's Matrix-Tree Theorem}

\textbf{Kirchhoff matrix.} The matrix $B$ of a graph $G$ of size $n \times n$ is defined as follows:
\[
B[i,j] = \begin{cases}
-1, & \text{if vertices numbered } i \text{ and } j \text{ are adjacent};\\[2pt]
0,  & \text{if } i \neq j \text{ and the vertices are not adjacent};\\[2pt]
\deg(v_i), & \text{if } i = j.
\end{cases}
\]

\textbf{Properties of the Kirchhoff matrix:}
\begin{enumerate}
  \item The sums of the elements in each row and each column of the matrix are equal to 0.
  \item The cofactors of all elements of the matrix are equal to one another.
  \item The determinant of the Kirchhoff matrix is zero.
\end{enumerate}

\textbf{Kirchhoff's matrix-tree theorem.} The number of spanning trees in a connected graph $G$ of order $n \geq 2$ is equal to the cofactor of any element of the Kirchhoff matrix $B$ of the graph $G$.

	\subsection{Boruvka's Algorithm}

Boruvka's algorithm is an algorithm for constructing a minimum spanning tree of a weighted connected undirected graph. It was proposed by Otakar Boruvka in 1926.

The algorithm works as follows:
\begin{enumerate}
  \item Initially, each vertex forms a separate component, that is, a forest of isolated vertices.
  \item For each component, the minimum-weight edge connecting it with another component is found.
  \item All found edges are added to the spanning tree, and the corresponding components are merged.
  \item Steps 2--3 are repeated until only one component remains.
\end{enumerate}

At each iteration, the number of components decreases at least by half, so there are at most $O(\log V)$ iterations. The total complexity is $O(E \log V)$.

	\subsection{Prufer Coding}

The Prufer encoding algorithm consists of the following steps. While the number of vertices is greater than two:
\begin{itemize}
  \item Choose the leaf $v$ with the smallest number.
  \item Add the number of the vertex adjacent to $v$ to the Prufer code.
  \item Remove vertex $v$ and its incident edge from the tree.
\end{itemize}

The obtained sequence is called the Prufer code of the given tree.

The decoding algorithm uses a Prufer code array of size $n-2$ and the array of all graph vertices $[1\ldots n]$. The following procedure is repeated $n-2$ times: the first element of the Prufer code is taken, and the smallest vertex not contained in the code is searched for in the vertex array. The found vertex and the current code element form an edge of the tree. These vertices are removed from the corresponding arrays. At the end, two vertices remain in the vertex array; they form the last edge of the tree. The complexity of the algorithms is $O(n)$.

	\subsection{Minimum Edge Cover}

An \textbf{edge cover} of a graph $G = (V, E)$ is a subset of edges $S \subseteq E$ such that every vertex of the graph is incident to at least one edge from $S$. A \textbf{minimum edge cover} is an edge cover of minimum cardinality.

The relation with maximum matching is expressed by Gallai's theorem: for a graph without isolated vertices,
\[
\alpha'(G) + \beta'(G) = |V(G)|,
\]
where $\alpha'(G)$ is the size of a maximum matching and $\beta'(G)$ is the size of a minimum edge cover.

The algorithm for finding a minimum edge cover is:
\begin{enumerate}
  \item Find a maximum matching $M$ in the graph.
  \item For each vertex not covered by matching $M$, add any edge incident to it to the cover.
\end{enumerate}

The complexity is determined by the complexity of finding a maximum matching.

	\subsection{Eulerian Graph}

An Eulerian graph is a graph that contains an Eulerian cycle, that is, a closed walk passing through all edges of the graph exactly once.

A graph is Eulerian if and only if it is connected and the degrees of all its vertices are even.

An Eulerian path is a path containing all edges. In this work, an Eulerian cycle is searched for. First, the program checks whether the graph is Eulerian. If it is not, the graph is modified so that the degree of every vertex becomes even.

To find an Eulerian cycle, an algorithm based on depth-first traversal is used:
\begin{itemize}
  \item Choose an arbitrary vertex of the graph and put it on the stack.
  \item While the stack is not empty, repeat the following actions:
  \begin{itemize}
    \item remove a vertex from the stack;
    \item if this vertex has unvisited adjacent vertices, choose one of them and put it on the stack;
    \item if the vertex has no unvisited adjacent vertices, add it to the Eulerian cycle list.
  \end{itemize}
  \item Return to the Eulerian cycle list in reverse order.
\end{itemize}

The complexity of finding an Eulerian cycle is estimated as $O(V+E)$.

	\subsection{Fundamental Cutset System}

A \textbf{cutset}, or cutting set, in a graph $G = (V, E)$ is a set of edges $C \subseteq E$ whose removal increases the number of connected components of the graph.

A \textbf{fundamental cutset} with respect to a spanning tree $T$ is the cutset corresponding to one tree edge $t \in E_T$: when $t$ is removed, the tree $T$ splits into two components $T_1$ and $T_2$; the fundamental cutset $C_t$ is the set of all edges of the original graph $G$ whose one endpoint lies in $T_1$ and the other endpoint lies in $T_2$.

The \textbf{fundamental cutset system} is constructed for a selected spanning tree $T$ and contains exactly $|V|-1$ cutsets, one for each tree edge.

All other cutsets of the graph can be obtained using the symmetric difference operation on fundamental cutsets:
\[
C = C_{t_1} \mathbin{\triangle} C_{t_2} \mathbin{\triangle} \cdots \mathbin{\triangle} C_{t_k},
\]
where $\triangle$ is the symmetric difference operation on sets.

\clearpage
\section{Implementation Features}

\subsection{Project Structure}

The program is implemented in C++17 and has a modular structure. Common data structures and generation functions are placed in the \texttt{shared} directory, the algorithms are divided by laboratory works, and the \texttt{lab\_combined} directory contains the combined console program with a menu. The full source code is given in Appendices A--E.

\begin{lstlisting}[caption={Project structure}, label={lst:project-structure}]
graph-theory-labs/
|-- CMakeLists.txt
|-- README.md
|-- shared/
|   |-- constants.h
|   |-- graph.h
|   |-- graph.cpp
|   |-- distribution.h
|   |-- distribution.cpp
|   |-- dag_generator.h
|   |-- dag_generator.cpp
|   |-- weight_matrix.h
|   \-- weight_matrix.cpp
|-- lab1/
|   |-- eccentricity.h
|   |-- eccentricity.cpp
|   |-- shimbell.h
|   |-- shimbell.cpp
|   |-- path_counter.h
|   \-- path_counter.cpp
|-- lab2/
|   |-- dfs.h
|   |-- dfs.cpp
|   |-- floyd_warshall.h
|   \-- floyd_warshall.cpp
|-- lab3/
|   |-- ford_fulkerson.h
|   |-- ford_fulkerson.cpp
|   |-- min_cost_flow.h
|   \-- min_cost_flow.cpp
|-- lab4/
|   |-- kirchhoff.h
|   |-- kirchhoff.cpp
|   |-- boruvka.h
|   |-- boruvka.cpp
|   |-- edge_cover.h
|   |-- edge_cover.cpp
|   |-- prufer.h
|   \-- prufer.cpp
|-- lab5/
|   |-- eulerian.h
|   |-- eulerian.cpp
|   |-- fundamental_cutsets.h
|   \-- fundamental_cutsets.cpp
\-- lab_combined/
    |-- CMakeLists.txt
    \-- main.cpp
\end{lstlisting}

The \texttt{shared} directory contains the graph class, the normal distribution generator, the graph generator, and the weight matrix generator. Directories \texttt{lab1}--\texttt{lab5} contain the algorithms of the corresponding laboratory works. The file \texttt{lab\_combined/main.cpp} connects all modules into a single user menu.

\subsection{Class Graph}

The \texttt{Graph} class is the basic structure for storing a directed or undirected graph. Its declaration is shown in Listing~\ref{lst:shared-graph-h}, and the method implementation is shown in Listing~\ref{lst:shared-graph-cpp}. The class stores the number of vertices \texttt{n}, the directedness flag \texttt{directed}, and the adjacency matrix \texttt{adj}.

The function \texttt{Graph(int n, bool directed)} creates an empty graph of the required type.

Input: the number of vertices \texttt{n} and the logical directedness flag \texttt{directed}.

Output: a graph object with an $n \times n$ adjacency matrix filled with zeros.

The function \texttt{addEdge} adds an edge or an arc between two vertices. If the graph is undirected, the adjacency matrix is filled symmetrically.

Input: vertex numbers \texttt{u} and \texttt{v}.

Output: the modified adjacency matrix of the graph.

The function \texttt{hasEdge} checks whether an edge or an arc exists.

Input: vertex numbers \texttt{u} and \texttt{v}.

Output: \texttt{true} if the edge exists; otherwise \texttt{false}.

The function \texttt{neighbors} returns all vertices adjacent to the specified vertex. For a directed graph it returns outgoing neighbors, and for an undirected graph it returns all neighboring vertices.

Input: vertex number \texttt{u}.

Output: a vector of neighboring vertex numbers.

The functions \texttt{printAdjMatrix} and \texttt{printFull} are used to print the graph to the console. The first one prints only the adjacency matrix, and the second one also prints the adjacency list, the edge list, and, if available, the weight matrix.

Input: a graph object and, for \texttt{printFull}, an optional weight matrix.

Output: a textual representation of the graph in the console.

\subsection{Distribution, Graph and Weight Matrix Generation}

Distribution generation functions are described in Listings~\ref{lst:shared-distribution-h} and~\ref{lst:shared-distribution-cpp}. Generation uses an approximate normal distribution formula through the sum of uniformly distributed random variables.

The function \texttt{normalSample} generates one random value approximately distributed according to the normal law.

Input: the number of uniform random variables \texttt{n} used in the approximation.

Output: a real number corresponding to the approximate normal distribution.

The function \texttt{sampleDegree} converts a normal-distribution value into an admissible vertex degree.

Input: the maximum admissible degree \texttt{maxDeg}.

Output: an integer from 1 to \texttt{maxDeg}.

The function \texttt{generateDegreeSequence} builds a degree sequence for the specified number of vertices. The sum of degrees is adjusted to $2(n-1)$, which corresponds to a connected acyclic undirected graph.

Input: the number of vertices \texttt{numVertices}.

Output: a vector of vertex degrees.

Graph generation is implemented in \texttt{dag\_generator.h} and \texttt{dag\_generator.cpp}, shown in Listings~\ref{lst:shared-dag-generator-h} and~\ref{lst:shared-dag-generator-cpp}. The main function \texttt{generateDAG} creates either a directed acyclic graph or an undirected graph obtained from the directed model.

Input: the number of vertices \texttt{n} and the directedness flag \texttt{directed}.

Output: a \texttt{Graph} object with a generated adjacency matrix.

For a directed graph, a random topological order is built first. Then a path ensuring connectivity is added, after which additional arcs are added only forward according to the topological order. Therefore, the graph remains acyclic. For an undirected graph, arc directions are ignored and the graph becomes symmetric.

Weight matrix generation is implemented in Listings~\ref{lst:shared-weight-matrix-h} and~\ref{lst:shared-weight-matrix-cpp}. The function \texttt{generateWeightMatrix} fills weights only for existing graph edges.

Input: graph \texttt{g} and weight generation mode \texttt{mode}.

Output: a weight matrix where absent edges are represented by \texttt{INF}.

Mode \texttt{POSITIVE} creates only positive weights, \texttt{NEGATIVE} creates only negative weights, and \texttt{MIXED} creates mixed weights. For an undirected graph, weights are written symmetrically.

\subsection{Laboratory Work No. 1}

The first laboratory work implements graph metric calculation, Shimbell's method, and route search. The code is located in \texttt{eccentricity.cpp}, \texttt{shimbell.cpp}, and \texttt{path\_counter.cpp}; the full texts are given in Listings~\ref{lst:lab1-eccentricity-cpp},~\ref{lst:lab1-shimbell-cpp}, and~\ref{lst:lab1-path-counter-cpp}.

The function \texttt{bfs} performs breadth-first search from a specified vertex and is used to calculate distances to all other vertices.

Input: graph \texttt{g}, start vertex \texttt{src}.

Output: a vector of distances from vertex \texttt{src}; value $-1$ means that a vertex is unreachable.

The function \texttt{computeMetrics} calculates eccentricities, radius, diameter, center, and diametral vertices of the graph. For each vertex it runs \texttt{bfs}, and then selects the maximum reachable distance from the obtained distances.

Input: graph \texttt{g}.

Output: the \texttt{GraphMetrics} structure containing the distance matrix, eccentricities, radius, diameter, center, and diametral vertices.

The function \texttt{shimbell} implements Shimbell's method for paths consisting of exactly $k$ edges. For minimum paths, tropical multiplication with minimum selection is used; for maximum paths, maximum selection is used.

Input: weight matrix \texttt{W}, number of edges \texttt{k}.

Output: a pair of matrices: the minimum-path matrix and the maximum-path matrix.

The function \texttt{findAllPaths} searches for all simple routes between two vertices. The implementation uses recursive depth-first search with a visited-vertex array to prevent a vertex from being repeated inside one route.

Input: graph \texttt{g}, start vertex \texttt{src}, destination vertex \texttt{dst}.

Output: a vector of routes, where each route is represented by a sequence of vertices.

The functions \texttt{pathExists} and \texttt{countPaths} are auxiliary wrappers. The first checks whether at least one route exists, and the second returns the number of found routes. The function \texttt{printPaths} prints routes to the console.

Input: a graph or an already found list of routes.

Output: a logical value, the number of routes, or textual output in the console.

\subsection{Laboratory Work No. 2}

The second laboratory work implements depth-first graph traversal and the Floyd-Warshall algorithm. The code is given in Listings~\ref{lst:lab2-dfs-cpp} and~\ref{lst:lab2-floyd-warshall-cpp}.

The function \texttt{dfs} performs depth-first traversal from a specified vertex. The implementation uses a stack, so the algorithm works iteratively. The number of iterations is also counted.

Input: graph \texttt{g}, start vertex \texttt{start}.

Output: traversal order, list of unreachable vertices, and the number of iterations printed to the console.

The function \texttt{floydWarshall} finds shortest paths between all pairs of vertices. Matrix \texttt{T} stores distances, and matrix \texttt{H} stores information for path reconstruction.

Input: graph \texttt{g}, weight matrix \texttt{W}.

Output: the \texttt{FloydResult} structure containing the distance matrix, the path reconstruction matrix, the number of iterations, and the negative-cycle flag.

The function \texttt{printFloydResult} prints the algorithm result and reconstructs the path between two vertices selected by the user. If a path does not exist, the program prints the corresponding message.

Input: the Floyd-Warshall result, the original weight matrix, the number of vertices, the start vertex, and the destination vertex.

Output: matrices \texttt{C}, \texttt{T}, \texttt{H}, the shortest path, and its length in the console.

\subsection{Laboratory Work No. 3}

The third laboratory work implements generation of capacity and cost matrices, the Ford-Fulkerson algorithm, and minimum-cost flow search. The code is located in Listings~\ref{lst:lab3-ford-fulkerson-cpp} and~\ref{lst:lab3-min-cost-flow-cpp}.

The function \texttt{generateCapacityMatrix} creates a capacity matrix for existing arcs of a directed graph.

Input: directed graph \texttt{g}.

Output: a capacity matrix where zero means that an arc is absent.

The function \texttt{generateCostMatrix} creates a matrix of the cost of passing one unit of flow through each arc.

Input: directed graph \texttt{g}.

Output: a cost matrix corresponding to the graph structure.

The function \texttt{fordFulkerson} implements the Ford-Fulkerson algorithm in the Edmonds-Karp variant, that is, an augmenting path is searched for using breadth-first search. At each step, the path capacity is found, and then the residual network is updated.

Input: number of vertices, capacity matrix, source \texttt{src}, sink \texttt{dst}.

Output: the \texttt{FFResult} structure containing the flow matrix and the maximum flow value.

The function \texttt{fordFulkersonSilent} performs the same algorithm but without printing intermediate paths. It is used inside the minimum-cost flow algorithm.

Input: number of vertices, capacity matrix, source and sink.

Output: the maximum flow value and the flow matrix without console output.

The function \texttt{minCostFlow} computes the minimum-cost flow for the value $\lfloor 2/3 \cdot \varphi_{\max} \rfloor$. First, the maximum flow is found, and then the residual network repeatedly selects shortest paths by cost found by the Floyd-Warshall algorithm.

Input: number of vertices, capacity matrix, cost matrix, source and sink.

Output: the \texttt{MCFResult} structure containing the maximum flow, required flow, achieved flow, total cost, and flow matrix.

The function \texttt{printMCFResult} prints the resulting minimum-cost flow and a verification calculation of the cost by arcs.

Input: the minimum-cost flow result, the capacity matrix, and the cost matrix.

Output: a table of arcs with positive flow and the total cost.

\subsection{Laboratory Work No. 4}

The fourth laboratory work implements Kirchhoff's matrix-tree theorem, Boruvka's algorithm, minimum edge cover, and Prufer code. The full code is given in Listings~\ref{lst:lab4-kirchhoff-cpp},~\ref{lst:lab4-boruvka-cpp},~\ref{lst:lab4-edge-cover-cpp}, and~\ref{lst:lab4-prufer-cpp}.

The function \texttt{countSpanningTrees} counts the number of spanning trees of a graph. For this purpose, the Kirchhoff matrix is built, the first row and the first column are removed, and the determinant of the resulting matrix is calculated.

Input: undirected graph \texttt{g}.

Output: the number of spanning trees of the graph.

The function \texttt{boruvka} builds a minimum spanning tree. At each iteration, the minimum outgoing edge is selected for each component, and then the components are merged.

Input: undirected graph \texttt{g}, weight matrix \texttt{W}.

Output: the \texttt{BoruvkaResult} structure containing the edges of the minimum spanning tree and its total weight.

The function \texttt{minEdgeCover} finds a minimum edge cover. First, a maximum matching is constructed; then incident edges are added for uncovered vertices.

Input: graph \texttt{g}, weight matrix, flag selecting the spanning tree or original graph, and the list of spanning tree edges.

Output: the \texttt{EdgeCoverResult} structure containing cover edges and the maximum matching size.

The function \texttt{pruferEncode} encodes the minimum spanning tree by the Prufer code. At each step, the leaf with the smallest number is removed, and its neighbor is written to the code. The weight of the removed edge is stored separately.

Input: list of spanning tree edges, number of vertices.

Output: the Prufer code and the list of edge weights in removal order.

The function \texttt{pruferDecode} restores a tree from the Prufer code and the saved weights. At each step, the minimum leaf absent from the current code is selected.

Input: Prufer code, list of weights, number of vertices.

Output: list of edges of the restored tree with weights.

\subsection{Laboratory Work No. 5}

The fifth laboratory work implements checking a graph for Eulerian property, building an Eulerian cycle, and constructing the fundamental cutset system. The code is given in Listings~\ref{lst:lab5-eulerian-cpp} and~\ref{lst:lab5-fundamental-cutsets-cpp}.

The function \texttt{buildEulerianCycle} checks graph connectivity and parity of all vertex degrees. If the graph is not Eulerian, the program attempts to modify it by adding or removing edges so that the degrees become even and connectivity is preserved.

Input: undirected graph \texttt{g}.

Output: the \texttt{EulerianResult} structure containing initial degrees, odd vertices, the list of changes, and the constructed Eulerian cycle.

The function \texttt{hierholzer} builds an Eulerian cycle after the graph has become Eulerian. The algorithm sequentially removes traversed edges and forms the cycle in reverse order.

Input: adjacency matrix of an Eulerian graph.

Output: sequence of vertices of the Eulerian cycle.

The function \texttt{buildFundamentalCutsets} constructs a fundamental cutset system with respect to the minimum spanning tree. For each spanning tree edge, the edge is temporarily removed; then two tree components are found, and all original graph edges connecting these components are selected.

Input: original graph \texttt{g}, list of minimum spanning tree edges.

Output: a vector of fundamental cutsets.

The function \texttt{printCutsetSymmetricDifference} computes the symmetric difference of selected fundamental cutsets. An edge that appears twice is removed from the result, while an edge that appears once remains.

Input: list of fundamental cutsets and numbers of selected cutsets.

Output: the set of edges obtained by the symmetric difference operation.

\subsection{Function main and Program Menu}

All laboratory works are combined in the file \texttt{lab\_combined/main.cpp}, whose full code is given in Listing~\ref{lst:lab-combined-main-cpp}. The program uses global variables to store the current graph, the weight matrix, capacity and cost matrices, the minimum spanning tree, and the Prufer encoding result.

The function \texttt{readInt} provides safe integer input. If the user enters invalid data, the input stream is cleared and the request is repeated.

Input: prompt string \texttt{prompt}.

Output: a correctly read integer.

The functions \texttt{requireGraph}, \texttt{requireWeight}, \texttt{requireFlowMatrices}, and \texttt{requireMST} check whether the required data has been prepared before an algorithm is launched.

Input: the current program state.

Output: an error message if the required data is missing.

Menu functions such as \texttt{menuGenerateGraph}, \texttt{menuShimbell}, \texttt{menuFloydWarshall}, \texttt{menuFordFulkerson}, \texttt{menuBoruvka}, and others read user parameters, check the correctness of the program state, and call the corresponding algorithms.

Input: the user's choice and additional values entered in the console.

Output: launch of the selected algorithm and printing of the result to the console.

The function \texttt{main} organizes the main program loop. Until the user selects the exit option, the menu is displayed, the command number is read, and the corresponding handler function is called.

Input: user commands from the console menu.

Output: sequential execution of laboratory algorithms until program termination.

Thus, the combined program makes it possible to perform the laboratory works in one interface: first the graph is generated, then the required matrices are built, and after that the user can launch algorithms from different sections of the work.

\subsection{Detailed Description of Functions}

Below is a more detailed description of the main functions of the program. For each function, its full name is given exactly as it is used in the source code, followed by its purpose, input data, and result. The source code is not duplicated in this section because it is fully provided in Appendices A--E; here the code is used as the basis for explaining the implementation. The function names in this subsection are hyperlinks to the corresponding source code listings.

\subsubsection*{Common Data Structures and Generation}

\noindent\textbf{Function} \coderef{Graph::Graph()}{lst:shared-graph-h}\\
The default constructor creates an empty graph object. It is needed for global variables and for cases where the graph is generated later through the program menu. Inside the constructor, the number of vertices is set to zero, and the graph is considered undirected until a new object is explicitly created.

\textbf{Input:} no input parameters.

\textbf{Output:} an empty \texttt{Graph} object for which the method \coderef{isEmpty()}{lst:shared-graph-h} returns \texttt{true}.

\noindent\textbf{Function} \coderef{Graph::Graph(int n, bool directed)}{lst:shared-graph-cpp}\\
The main constructor of the \texttt{Graph} class. It creates an adjacency matrix of size $n \times n$ and fills it with zeros. The parameter \texttt{directed} determines whether the graph is directed. Since the graph is stored as an adjacency matrix, checking whether an edge exists takes constant time.

\textbf{Input:} the number of vertices \texttt{n}; logical parameter \texttt{directed}, equal to \texttt{true} for a directed graph and \texttt{false} for an undirected graph.

\textbf{Output:} a graph object with a prepared adjacency matrix and a stored graph type.

\noindent\textbf{Function} \coderef{Graph::addEdge(int u, int v)}{lst:shared-graph-cpp}\\
The method adds an edge between vertices \texttt{u} and \texttt{v}. If the graph is directed, only the element \texttt{adj[u][v]} is set in the adjacency matrix. If the graph is undirected, the symmetric element \texttt{adj[v][u]} is also set, so one edge is stored in two matrix cells.

\textbf{Input:} vertex numbers \texttt{u} and \texttt{v} in internal zero-based numbering.

\textbf{Output:} the modified adjacency matrix of the current \texttt{Graph} object.

\noindent\textbf{Function} \coderef{Graph::hasEdge(int u, int v) const}{lst:shared-graph-cpp}\\
The method checks whether an edge or arc from vertex \texttt{u} to vertex \texttt{v} exists. The implementation accesses one cell of the adjacency matrix and compares it with one.

\textbf{Input:} vertex numbers \texttt{u} and \texttt{v}.

\textbf{Output:} \texttt{true} if \texttt{adj[u][v] == 1}; otherwise \texttt{false}.

\noindent\textbf{Function} \coderef{Graph::neighbors(int u) const}{lst:shared-graph-cpp}\\
The method forms a list of all vertices reachable from vertex \texttt{u} in one step. For a directed graph, these are outgoing neighbors; for an undirected graph, these are all adjacent vertices because the adjacency matrix is symmetric.

\textbf{Input:} vertex number \texttt{u}.

\textbf{Output:} a \lstinline|std::vector<int>| vector with neighboring vertex numbers.

\noindent\textbf{Function} \coderef{Graph::printAdjMatrix() const}{lst:shared-graph-cpp}\\
The method prints the adjacency matrix to the console. During output, vertices are converted from internal numbering $0,1,\ldots,n-1$ to user numbering $1,2,\ldots,n$. For an absent edge, \texttt{inf} is printed; for an existing edge, \texttt{1} is printed.

\textbf{Input:} the current \texttt{Graph} object.

\textbf{Output:} the adjacency matrix table in the console.

\noindent\textbf{Function} \coderef{Graph::printFull(const std::vector<std::vector<double>>* weightMatrix) const}{lst:shared-graph-cpp}\\
The method prints an extended representation of the graph: adjacency list, edge list, adjacency matrix, and, if passed, the weight matrix. The pointer to the weight matrix is optional, so the same function is used both before and after weight generation.

\textbf{Input:} the current graph and an optional pointer \texttt{weightMatrix} to the weight matrix.

\textbf{Output:} a complete textual description of the graph in the console.

\noindent\textbf{Function} \coderef{Graph::isEmpty() const}{lst:shared-graph-h}\\
The method is used to quickly check whether a graph has already been created. It returns true if the number of vertices is zero.

\textbf{Input:} the current \texttt{Graph} object.

\textbf{Output:} a logical value showing whether the graph is empty.

\noindent\textbf{Function} \coderef{double normalSample(int n)}{lst:shared-distribution-cpp}\\
The function generates one value approximately distributed according to the normal law. It uses the sum of \texttt{n} uniformly distributed random variables and normalization according to the formula shown in the source file comments. This approach is convenient because it is simple and does not require a separate complex normal-distribution generator.

\textbf{Input:} number \texttt{n}, which sets the number of uniform random variables for approximation.

\textbf{Output:} a \texttt{double} value imitating a sample from the normal distribution.

\noindent\textbf{Function} \coderef{int sampleDegree(int maxDeg)}{lst:shared-distribution-cpp}\\
The function converts a random normal-distribution value into an integer vertex degree. The obtained number is limited to the range from \texttt{1} to \texttt{maxDeg}, so every vertex receives an admissible positive degree.

\textbf{Input:} maximum admissible degree \texttt{maxDeg}.

\textbf{Output:} an integer representing a vertex degree.

\noindent\textbf{Function} \coderef{std::vector<int> generateDegreeSequence(int numVertices)}{lst:shared-distribution-cpp}\\
The function creates a degree sequence for the specified number of vertices. First, a degree is generated for each vertex; then the degree sum is adjusted so that it corresponds to the required graph generation model. This sequence is later used when constructing a directed acyclic graph.

\textbf{Input:} number of vertices \texttt{numVertices}.

\textbf{Output:} a vector of integers where each element corresponds to the degree of one vertex.

\noindent\textbf{Function} \coderef{Graph generateDAG(int n, bool directed)}{lst:shared-dag-generator-cpp}\\
The function is the main graph generation entry point. If \texttt{directed} is \texttt{true}, a directed acyclic graph is built. A random topological order is created, a path for connectivity is added, and additional arcs are placed only forward in this order. If \texttt{directed} is \texttt{false}, a directed model is built first, then edge directions are ignored, producing an undirected graph.

\textbf{Input:} number of vertices \texttt{n}; directedness flag \texttt{directed}.

\textbf{Output:} the generated \texttt{Graph} object.

\noindent\textbf{Function} \coderef{Matrix generateWeightMatrix(const Graph& g, WeightMode mode)}{lst:shared-weight-matrix-cpp}\\
The function creates a weight matrix for an already generated graph. A weight is assigned only where an edge exists in the graph. For absent edges, a special infinity value is used, which is important for Shimbell's method and the Floyd-Warshall algorithm. Mode \texttt{POSITIVE} creates positive weights, \texttt{NEGATIVE} creates negative weights, and \texttt{MIXED} creates mixed weights.

\textbf{Input:} graph \texttt{g}; weight generation mode \texttt{mode}.

\textbf{Output:} a \texttt{Matrix} of size $n \times n$ with weights of existing edges.

\subsubsection*{Functions of Laboratory Work No. 1}

\noindent\textbf{Function} \coderef{std::vector<int> bfs(const Graph& g, int src)}{lst:lab1-eccentricity-cpp}\\
The function performs breadth-first search from vertex \texttt{src}. It is used as an auxiliary part of calculating graph metrics. The search follows the neighbor list obtained from the adjacency matrix. If a vertex is reachable, the shortest distance in the number of edges is written for it; if a vertex is unreachable, the value remains \texttt{-1}.

\textbf{Input:} graph \texttt{g}; start vertex \texttt{src}.

\textbf{Output:} a vector of distances from vertex \texttt{src} to all other vertices.

\noindent\textbf{Function} \coderef{GraphMetrics computeMetrics(const Graph& g)}{lst:lab1-eccentricity-cpp}\\
The function calculates the main metric characteristics of a graph. For each vertex, \coderef{bfs}{lst:lab1-eccentricity-cpp} is called, after which a distance matrix is formed. The eccentricity of a vertex is the maximum distance from it to reachable vertices. Then the radius, diameter, center vertices, and diametral vertices are selected from the eccentricities.

\textbf{Input:} graph \texttt{g}.

\textbf{Output:} the \texttt{GraphMetrics} structure containing the distance matrix, eccentricities, radius, diameter, center, and diametral vertices.

\noindent\textbf{Function} \coderef{void printMetrics(const GraphMetrics& m)}{lst:lab1-eccentricity-cpp}\\
The function prints the results obtained in \coderef{computeMetrics}{lst:lab1-eccentricity-cpp}: the distance matrix, eccentricities, radius, diameter, center, and diametral vertices.

\textbf{Input:} the \texttt{GraphMetrics} structure.

\textbf{Output:} formatted output of metric characteristics in the console.

\noindent\textbf{Function} \coderef{std::pair<Matrix, Matrix> shimbell(const Matrix& W, int k)}{lst:lab1-shimbell-cpp}\\
The function implements Shimbell's method for finding paths consisting of exactly \texttt{k} edges. Internally, tropical matrix multiplication is used: for minimum paths, the minimum sum of weights is selected; for maximum paths, the maximum is selected. For \texttt{k = 0}, an analogue of the identity matrix is returned, where the path from a vertex to itself has length zero.

\textbf{Input:} weight matrix \texttt{W}; number of edges in the path \texttt{k}.

\textbf{Output:} a pair of matrices: the first contains minimum path weights, and the second contains maximum path weights for exactly \texttt{k} steps.

\noindent\textbf{Function} \coderef{void printMatrix(const Matrix& M)}{lst:lab1-shimbell-cpp}\\
The function prints a weight matrix to the screen. Large positive values are displayed as \texttt{INF}, and large negative values are displayed as \texttt{-INF}.

\textbf{Input:} matrix \texttt{M}.

\textbf{Output:} a table of matrix values in the console.

\noindent\textbf{Function} \coderef{std::vector<std::vector<int>> findAllPaths(const Graph& g, int src, int dst)}{lst:lab1-path-counter-cpp}\\
The function finds all simple routes between vertices \texttt{src} and \texttt{dst}. Recursive depth-first search is used. The visited-vertex array prevents the same vertex from being repeated in one route.

\textbf{Input:} graph \texttt{g}; start vertex \texttt{src}; destination vertex \texttt{dst}.

\textbf{Output:} a list of routes, where each route is represented by a vector of vertex numbers.

\noindent\textbf{Function} \coderef{bool pathExists(const Graph& g, int src, int dst)}{lst:lab1-path-counter-cpp}\\
The function checks whether at least one path exists between the selected vertices.

\textbf{Input:} graph \texttt{g}; start vertex \texttt{src}; destination vertex \texttt{dst}.

\textbf{Output:} \texttt{true} if a path exists; otherwise \texttt{false}.

\noindent\textbf{Function} \coderef{int countPaths(const Graph& g, int src, int dst)}{lst:lab1-path-counter-cpp}\\
The function counts the number of simple routes between two vertices.

\textbf{Input:} graph \texttt{g}; start vertex \texttt{src}; destination vertex \texttt{dst}.

\textbf{Output:} an integer equal to the number of found routes.

\noindent\textbf{Function} \coderef{void printPaths(const std::vector<std::vector<int>>& paths, int src, int dst)}{lst:lab1-path-counter-cpp}\\
The function prints the found routes using user vertex numbering. If the list is empty, it reports that there are no routes between the selected vertices.

\textbf{Input:} list of routes \texttt{paths}; start vertex \texttt{src}; destination vertex \texttt{dst}.

\textbf{Output:} a textual list of routes in the console.

\subsubsection*{Functions of Laboratory Work No. 2}

\noindent\textbf{Function} \coderef{void dfs(const Graph& g, int start)}{lst:lab2-dfs-cpp}\\
The function performs depth-first traversal of the graph starting from vertex \texttt{start}. The implementation uses a stack, records the visiting order, counts iterations, and prints unreachable vertices.

\textbf{Input:} graph \texttt{g}; start vertex \texttt{start}.

\textbf{Output:} traversal order, list of unreachable vertices, and number of iterations in the console.

\noindent\textbf{Function} \coderef{FloydResult floydWarshall(const Graph& g, const std::vector<std::vector<double>>& W)}{lst:lab2-floyd-warshall-cpp}\\
The function implements the Floyd-Warshall algorithm for finding shortest paths between all pairs of vertices. Matrix \texttt{T} stores current distances, and matrix \texttt{H} is used to reconstruct paths. The diagonal elements are checked to detect negative cycles.

\textbf{Input:} graph \texttt{g}; weight matrix \texttt{W}, where the diagonal is zero for this algorithm.

\textbf{Output:} the \texttt{FloydResult} structure with the distance matrix, path reconstruction matrix, number of iterations, and negative-cycle flag.

\noindent\textbf{Function} \coderef{void printFloydResult(const FloydResult& res, const std::vector<std::vector<double>>& W, int n, int src, int dst)}{lst:lab2-floyd-warshall-cpp}\\
The function prints the original weight matrix, the final shortest-distance matrix, the path reconstruction matrix, and the path between vertices selected by the user.

\textbf{Input:} result \texttt{res}; weight matrix \texttt{W}; number of vertices \texttt{n}; start vertex \texttt{src}; destination vertex \texttt{dst}.

\textbf{Output:} detailed output of tables and the reconstructed shortest path in the console.

\subsubsection*{Functions of Laboratory Work No. 3}

\noindent\textbf{Function} \coderef{std::vector<std::vector<int>> generateCapacityMatrix(const Graph& g)}{lst:lab3-ford-fulkerson-cpp}\\
The function creates a capacity matrix for a directed graph. For each existing arc, an integer capacity is generated, and absent arcs remain zero.

\textbf{Input:} directed graph \texttt{g}.

\textbf{Output:} a capacity matrix where a positive number means that an arc exists and gives its capacity.

\noindent\textbf{Function} \coderef{std::vector<std::vector<int>> generateCostMatrix(const Graph& g)}{lst:lab3-ford-fulkerson-cpp}\\
The function creates a cost matrix for the same arcs that are present in the graph.

\textbf{Input:} directed graph \texttt{g}.

\textbf{Output:} a cost matrix where zero for an absent arc means that flow along it is not allowed.

\noindent\textbf{Function} \coderef{FFResult fordFulkerson(int n, const std::vector<std::vector<int>>& cap, int src, int dst)}{lst:lab3-ford-fulkerson-cpp}\\
The function implements the Ford-Fulkerson algorithm in the Edmonds-Karp variant. Every augmenting path is found by BFS, then the bottleneck capacity is used to increase the flow and update reverse arcs.

\textbf{Input:} number of vertices \texttt{n}; capacity matrix \texttt{cap}; source \texttt{src}; sink \texttt{dst}.

\textbf{Output:} the \texttt{FFResult} structure containing the final flow matrix, maximum flow value, source, and sink.

\noindent\textbf{Function} \coderef{FFResult fordFulkersonSilent(int n, const std::vector<std::vector<int>>& cap, int src, int dst)}{lst:lab3-ford-fulkerson-cpp}\\
The function performs the same maximum-flow calculation as \coderef{fordFulkerson}{lst:lab3-ford-fulkerson-cpp}, but it does not print intermediate augmenting paths.

\textbf{Input:} number of vertices \texttt{n}; capacity matrix \texttt{cap}; source \texttt{src}; sink \texttt{dst}.

\textbf{Output:} the \texttt{FFResult} structure without intermediate console output.

\noindent\textbf{Function} \coderef{MCFResult minCostFlow(int n, const std::vector<std::vector<int>>& cap, const std::vector<std::vector<int>>& cost, int src, int dst)}{lst:lab3-min-cost-flow-cpp}\\
The function computes the minimum-cost flow. First, the maximum flow is found using \coderef{fordFulkersonSilent}{lst:lab3-ford-fulkerson-cpp}; then the target flow is set to $\lfloor 2/3 \cdot \varphi_{max} \rfloor$. The flow is accumulated by adding shortest paths by cost in the residual network.

\textbf{Input:} number of vertices \texttt{n}; capacity matrix \texttt{cap}; cost matrix \texttt{cost}; source \texttt{src}; sink \texttt{dst}.

\textbf{Output:} the \texttt{MCFResult} structure with maximum flow, target flow, achieved flow, total cost, and flow matrix.

\subsubsection*{Functions of Laboratory Work No. 4}

\noindent\textbf{Function} \coderef{long long countSpanningTrees(const Graph& g)}{lst:lab4-kirchhoff-cpp}\\
The function counts the number of spanning trees using Kirchhoff's matrix-tree theorem by constructing the Laplacian matrix, deleting one row and one column, and computing the determinant.

\textbf{Input:} undirected graph \texttt{g}.

\textbf{Output:} the number of spanning trees as \texttt{long long}.

\noindent\textbf{Function} \coderef{BoruvkaResult boruvka(const Graph& g, const std::vector<std::vector<double>>& W)}{lst:lab4-boruvka-cpp}\\
The function implements Boruvka's algorithm. At each iteration, the cheapest outgoing edge is selected for each component; selected edges are added to the spanning tree and components are merged.

\textbf{Input:} undirected graph \texttt{g}; weight matrix \texttt{W}.

\textbf{Output:} the \texttt{BoruvkaResult} structure containing minimum spanning tree edges and total weight.

\noindent\textbf{Function} \coderef{EdgeCoverResult minEdgeCover(const Graph& g, const std::vector<std::vector<double>>& W, bool useMST, const std::vector<MSTEdge>& mstEdges)}{lst:lab4-edge-cover-cpp}\\
The function builds a minimum edge cover either on the minimum spanning tree or on the full graph. It first finds a maximum matching and then adds incident edges for uncovered vertices.

\textbf{Input:} graph \texttt{g}; weight matrix \texttt{W}; flag \texttt{useMST}; list of spanning tree edges \texttt{mstEdges}.

\textbf{Output:} the \texttt{EdgeCoverResult} structure containing cover edges and matching size.

\noindent\textbf{Function} \coderef{PruferResult pruferEncode(const std::vector<MSTEdge>& mstEdges, int n)}{lst:lab4-prufer-cpp}\\
The function encodes a spanning tree using the Prufer code. It repeatedly removes the smallest leaf, records its neighbor, and saves the removed edge weight.

\textbf{Input:} list of spanning tree edges \texttt{mstEdges}; number of vertices \texttt{n}.

\textbf{Output:} the \texttt{PruferResult} structure containing the Prufer code and saved weights.

\noindent\textbf{Function} \coderef{std::vector<MSTEdge> pruferDecode(const std::vector<int>& code, const std::vector<double>& weights, int n)}{lst:lab4-prufer-cpp}\\
The function restores a tree from a Prufer code and the saved weights by repeatedly connecting the smallest leaf absent from the current code with the first code element.

\textbf{Input:} Prufer code \texttt{code}; list of weights \texttt{weights}; number of vertices \texttt{n}.

\textbf{Output:} list of edges of the restored tree.

\subsubsection*{Functions of Laboratory Work No. 5}

\noindent\textbf{Function} \coderef{EulerianResult buildEulerianCycle(const Graph& g)}{lst:lab5-eulerian-cpp}\\
The function checks whether an undirected graph is Eulerian and builds an Eulerian cycle. If the conditions are not satisfied, it tries to modify a local copy of the graph without changing the original object.

\textbf{Input:} undirected graph \texttt{g}.

\textbf{Output:} the \texttt{EulerianResult} structure containing original degrees, odd vertices, performed changes, and the found Eulerian cycle.

\noindent\textbf{Function} \coderef{std::vector<FundamentalCutset> buildFundamentalCutsets(const Graph& g, const std::vector<MSTEdge>& mstEdges)}{lst:lab5-fundamental-cutsets-cpp}\\
The function builds the fundamental cutset system with respect to the minimum spanning tree. For every tree edge, the edge is removed, two components are found, and all original graph edges connecting these components are collected.

\textbf{Input:} original graph \texttt{g}; list of minimum spanning tree edges \texttt{mstEdges}.

\textbf{Output:} a vector of \texttt{FundamentalCutset} structures.

\noindent\textbf{Function} \coderef{void printCutsetSymmetricDifference(const std::vector<FundamentalCutset>& cutsets, const std::vector<int>& selectedIndices)}{lst:lab5-fundamental-cutsets-cpp}\\
The function computes the symmetric difference of selected fundamental cutsets, keeping edges that occur an odd number of times.

\textbf{Input:} list of cutsets \texttt{cutsets}; numbers of selected cutsets \texttt{selectedIndices}.

\textbf{Output:} the set of edges obtained as the symmetric difference, printed to the console.

\subsubsection*{Functions of the Combined Menu}

\noindent\textbf{Function} \coderef{static inline int D(int v)}{lst:lab-combined-main-cpp}\\
The function converts internal zero-based vertex numbering to user numbering starting from one.

\textbf{Input:} vertex number \texttt{v} in internal numbering.

\textbf{Output:} vertex number \texttt{v + 1} for console output.

\noindent\textbf{Function} \coderef{static int readInt(const std::string& prompt)}{lst:lab-combined-main-cpp}\\
The function safely reads an integer, repeats the request on invalid input, and clears the input stream.

\textbf{Input:} prompt string \texttt{prompt}.

\textbf{Output:} correctly read integer.

\noindent\textbf{Functions} \coderef{requireGraph()}{lst:lab-combined-main-cpp}, \coderef{requireWeight()}{lst:lab-combined-main-cpp}, \coderef{requireFlowMatrices()}{lst:lab-combined-main-cpp}, \coderef{requireMST()}{lst:lab-combined-main-cpp}\\
These functions print error messages if the user tries to launch an algorithm without the required preliminary data.

\textbf{Input:} current state of global program variables.

\textbf{Output:} diagnostic message in the console.

\noindent\textbf{Functions} \coderef{void menuGenerateGraph()}{lst:lab-combined-main-cpp}, \coderef{void menuGenerateWeightMatrix()}{lst:lab-combined-main-cpp}, \coderef{menuDFS()}{lst:lab-combined-main-cpp}, \coderef{menuFloydWarshall()}{lst:lab-combined-main-cpp}, \coderef{menuFordFulkerson()}{lst:lab-combined-main-cpp}, \coderef{menuBoruvka()}{lst:lab-combined-main-cpp}, \coderef{menuEulerianCycle()}{lst:lab-combined-main-cpp}\\
These menu functions read user parameters, check that the current program state is valid, and call the corresponding algorithmic functions.

\textbf{Input:} current graph, matrices, and parameters entered by the user.

\textbf{Output:} execution of the selected algorithm and result output in the console.

\noindent\textbf{Function} \coderef{int main()}{lst:lab-combined-main-cpp}\\
The main function organizes the general application loop. Until the user selects item \texttt{0}, the program prints the menu, reads the command number through \coderef{readInt}{lst:lab-combined-main-cpp}, and calls the corresponding menu function using the \texttt{switch} operator.

\textbf{Input:} user commands from the console.

\textbf{Output:} execution of the selected algorithms and program termination after the exit item is selected.

\clearpage
\section{Program Results}

\subsection{Tasks 1--11: Directed Graph with 5 Vertices}

After the program starts, the main menu is displayed, containing items for all five laboratory works. For tasks 1--11, a directed graph with 5 vertices was generated, because the first three laboratory works use a directed acyclic graph and algorithms for routes, shortest paths, and flows. The main program menu is shown in Fig.~\ref{fig:res-ui}.

\begin{figure}[H]
	\centering
	\includegraphics[width=0.95\textwidth,height=0.78\textheight,keepaspectratio]{\detokenize{1 to 11 with directed graph vertices 5/ui.png}}
	\caption{Main menu of the combined program}
	\label{fig:res-ui}
\end{figure}

When item \texttt{1} is selected, the user enters the number of vertices and the graph type. For the first part of the results, a directed graph was selected. The program builds the adjacency matrix and reports that a weight matrix must be generated for further algorithms; see Fig.~\ref{fig:res-generate-directed}.

\begin{figure}[H]
	\centering
	\includegraphics[width=0.95\textwidth,height=0.78\textheight,keepaspectratio]{\detokenize{1 to 11 with directed graph vertices 5/1.generate graph directed Vertices 5.png}}
	\caption{Generation of a directed graph with 5 vertices}
	\label{fig:res-generate-directed}
\end{figure}

When item \texttt{2} is selected, the full graph representation is printed: adjacency list, arc list, and adjacency matrix. This makes it possible to check the structure of the generated directed graph before launching algorithms; see Fig.~\ref{fig:res-show-directed}.

\begin{figure}[H]
	\centering
	\includegraphics[width=0.95\textwidth,height=0.78\textheight,keepaspectratio]{\detokenize{1 to 11 with directed graph vertices 5/2. show graph V5.png}}
	\caption{Output of the adjacency list, arc list, and adjacency matrix}
	\label{fig:res-show-directed}
\end{figure}

When item \texttt{3} is selected, the user chooses the weight generation mode: positive, negative, or mixed. After that, the program prints the weight matrix, where INF means that an arc is absent; see Fig.~\ref{fig:res-weight-directed}.

\begin{figure}[H]
	\centering
	\includegraphics[width=0.95\textwidth,height=0.78\textheight,keepaspectratio]{\detokenize{1 to 11 with directed graph vertices 5/3. generate weight matrix.png}}
	\caption{Weight matrix generation for the directed graph}
	\label{fig:res-weight-directed}
\end{figure}

When item \texttt{4} is selected, distances between vertices, eccentricities, radius, diameter, graph center, and diametral vertices are calculated. The algorithm result is represented as a distance matrix and a final table; see Fig.~\ref{fig:res-eccentricity}.

\begin{figure}[H]
	\centering
	\includegraphics[width=0.95\textwidth,height=0.78\textheight,keepaspectratio]{\detokenize{1 to 11 with directed graph vertices 5/4. Eccentricities, center, diametral vertices.png}}
	\caption{Graph eccentricities, center, and diametral vertices}
	\label{fig:res-eccentricity}
\end{figure}

When item \texttt{5} is selected, the user specifies the number of steps $k$, after which the program applies Shimbell's method. The result contains matrices of minimum and maximum paths consisting of exactly the specified number of arcs; see Fig.~\ref{fig:res-shimbell}.

\begin{figure}[H]
	\centering
	\includegraphics[width=0.95\textwidth,height=0.78\textheight,keepaspectratio]{\detokenize{1 to 11 with directed graph vertices 5/5. shimbell's method.png}}
	\caption{Shimbell's method for the directed weighted graph}
	\label{fig:res-shimbell}
\end{figure}

When item \texttt{6} is selected, the user enters the start and destination vertices. The program determines whether a route exists between them, counts the number of simple routes, and prints the found routes; see Fig.~\ref{fig:res-routes}.

\begin{figure}[H]
	\centering
	\includegraphics[width=0.95\textwidth,height=0.78\textheight,keepaspectratio]{\detokenize{1 to 11 with directed graph vertices 5/6. find all routes between two vertices.png}}
	\caption{Search for all routes between two vertices}
	\label{fig:res-routes}
\end{figure}

When item \texttt{7} is selected, depth-first graph traversal is performed from the specified vertex. The result contains the traversal order, the number of visited vertices, and the number of algorithm iterations; see Fig.~\ref{fig:res-dfs}.

\begin{figure}[H]
	\centering
	\includegraphics[width=0.95\textwidth,height=0.72\textheight,keepaspectratio]{\detokenize{1 to 11 with directed graph vertices 5/7. DFS traversal.png}}
	\caption{Depth-first graph traversal}
	\label{fig:res-dfs}
\end{figure}

When item \texttt{8} is selected, the user enters the start and destination vertices for shortest-path search. The Floyd-Warshall algorithm prints the original weight matrix, shortest-distance matrix, path reconstruction matrix, the path itself, and the number of iterations; see Fig.~\ref{fig:res-floyd}.

\begin{figure}[H]
	\centering
	\includegraphics[width=0.95\textwidth,height=0.78\textheight,keepaspectratio]{\detokenize{1 to 11 with directed graph vertices 5/8. Shortest path Floyd-warshall.png}}
	\caption{Shortest path search using the Floyd-Warshall algorithm}
	\label{fig:res-floyd}
\end{figure}

When item \texttt{9} is selected, capacity and cost matrices are generated based on the directed graph. In addition, the program prints incoming and outgoing vertex degrees so that the user can choose the source and sink more easily; see Fig.~\ref{fig:res-flow-matrices}.

\begin{figure}[H]
	\centering
	\includegraphics[width=0.95\textwidth,height=0.78\textheight,keepaspectratio]{\detokenize{1 to 11 with directed graph vertices 5/9. generate capacity and cost matrices.png}}
	\caption{Generation of capacity and cost matrices}
	\label{fig:res-flow-matrices}
\end{figure}

When item \texttt{10} is selected, the program finds the maximum flow using the Ford-Fulkerson algorithm. The screen displays augmenting paths, the maximum-flow value, and flow distribution by arcs; see Fig.~\ref{fig:res-ford-fulkerson}.

\begin{figure}[H]
	\centering
	\includegraphics[width=0.95\textwidth,height=0.78\textheight,keepaspectratio]{\detokenize{1 to 11 with directed graph vertices 5/10. Max flow ford fulkerson.png}}
	\caption{Finding the maximum flow using the Ford-Fulkerson algorithm}
	\label{fig:res-ford-fulkerson}
\end{figure}

When item \texttt{11} is selected, the minimum-cost flow is computed. The required flow value is $\lfloor 2/3 \cdot \varphi_{\max} \rfloor$. The program prints the shortest paths by cost, the achieved flow, and the total cost; see Fig.~\ref{fig:res-min-cost-flow}.

\begin{figure}[H]
	\centering
	\includegraphics[width=0.95\textwidth,height=0.78\textheight,keepaspectratio]{\detokenize{1 to 11 with directed graph vertices 5/11. Min-cost flow.png}}
	\caption{Minimum-cost flow calculation}
	\label{fig:res-min-cost-flow}
\end{figure}

\clearpage
\subsection{Tasks 12--17: Undirected Graph with 7 Vertices}

For tasks 12--17, an undirected graph with 7 vertices was generated, because the algorithms of the fourth and fifth laboratory works operate on an undirected graph. After generation, the adjacency matrix is printed; see Fig.~\ref{fig:res-generate-undirected}.

\begin{figure}[H]
	\centering
	\includegraphics[width=0.95\textwidth,height=0.78\textheight,keepaspectratio]{\detokenize{12 to 17 with undirected graph vertices 7/1. generate graph for 12 to 17 lab 4 and 5.png}}
	\caption{Generation of an undirected graph with 7 vertices}
	\label{fig:res-generate-undirected}
\end{figure}

When item \texttt{12} is selected, the program builds the Kirchhoff matrix, removes the first row and first column to obtain a cofactor, and calculates the determinant. The obtained value is the number of spanning trees of the graph; see Fig.~\ref{fig:res-kirchhoff}.

\begin{figure}[H]
	\centering
	\includegraphics[width=0.95\textwidth,height=0.78\textheight,keepaspectratio]{\detokenize{12 to 17 with undirected graph vertices 7/12. kirchhoff count spanning trees.png}}
	\caption{Counting spanning trees using Kirchhoff's theorem}
	\label{fig:res-kirchhoff}
\end{figure}

When item \texttt{13} is selected, a minimum spanning tree is built using Boruvka's algorithm. The program prints connected components, added safe edges, and the final weight of the obtained spanning tree step by step; see Fig.~\ref{fig:res-boruvka}.

\begin{figure}[H]
	\centering
	\includegraphics[width=0.95\textwidth,height=0.78\textheight,keepaspectratio]{\detokenize{12 to 17 with undirected graph vertices 7/13 Boruvka minimun spanning tree.png}}
	\caption{Building a minimum spanning tree using Boruvka's algorithm}
	\label{fig:res-boruvka}
\end{figure}

When item \texttt{14} is selected, the user chooses whether to search for a minimum edge cover on the original graph or on the constructed spanning tree. The result contains the size of the maximum matching and the edges of the minimum cover; see Fig.~\ref{fig:res-edge-cover}.

\begin{figure}[H]
	\centering
	\includegraphics[width=0.95\textwidth,height=0.78\textheight,keepaspectratio]{\detokenize{12 to 17 with undirected graph vertices 7/14. Minimum edge cover.png}}
	\caption{Finding a minimum edge cover}
	\label{fig:res-edge-cover}
\end{figure}

When item \texttt{15} is selected, the constructed minimum spanning tree is encoded by the Prufer code with weight preservation, and then reverse decoding is performed. The screen displays removed leaves, the Prufer code, edge weights, and the restored tree; see Fig.~\ref{fig:res-prufer}.

\begin{figure}[H]
	\centering
	\includegraphics[width=0.95\textwidth,height=0.78\textheight,keepaspectratio]{\detokenize{12 to 17 with undirected graph vertices 7/15. prufer encode decode.png}}
	\caption{Encoding and decoding the spanning tree using the Prufer code}
	\label{fig:res-prufer}
\end{figure}

When item \texttt{16} is selected, the program checks whether the graph is Eulerian. If the graph does not satisfy the even-degree condition, edge modification is performed; after that, an Eulerian cycle is built; see Fig.~\ref{fig:res-eulerian}.

\begin{figure}[H]
	\centering
	\includegraphics[width=0.95\textwidth,height=0.78\textheight,keepaspectratio]{\detokenize{12 to 17 with undirected graph vertices 7/16. Eulerian check modify .png}}
	\caption{Checking Eulerian property, modification, and Eulerian cycle construction}
	\label{fig:res-eulerian}
\end{figure}

When item \texttt{17} is selected, the program builds a minimum spanning tree using Boruvka's algorithm and then obtains a fundamental cutset for every edge of the spanning tree. After printing the fundamental system, the user selects cutset numbers, and the program computes their symmetric difference; see Fig.~\ref{fig:res-cutsets-1} and Fig.~\ref{fig:res-cutsets-2}.

\begin{figure}[H]
	\centering
	\includegraphics[width=0.95\textwidth,height=0.78\textheight,keepaspectratio]{\detokenize{12 to 17 with undirected graph vertices 7/17.1 fundamental cutsets form MST.png}}
	\caption{Construction of the fundamental cutset system}
	\label{fig:res-cutsets-1}
\end{figure}

\begin{figure}[H]
	\centering
	\includegraphics[width=0.95\textwidth,height=0.78\textheight,keepaspectratio]{\detokenize{12 to 17 with undirected graph vertices 7/17.2 fundamental cutsets form MST.png}}
	\caption{Symmetric difference of selected fundamental cutsets}
	\label{fig:res-cutsets-2}
\end{figure}

\clearpage
\subsection{Error Handling}

The program checks the correctness of user actions. If the user tries to call an algorithm before generating a graph or a weight matrix, the program prints a message indicating the required menu item; see Fig.~\ref{fig:res-error-1}.

\begin{figure}[H]
	\centering
	\includegraphics[width=0.9\textwidth,height=0.45\textheight,keepaspectratio]{\detokenize{error handling/error handle1.png}}
	\caption{Error message when required data is missing}
	\label{fig:res-error-1}
\end{figure}

For algorithms of the fourth and fifth laboratory works, the graph type is additionally checked. If the user launches an algorithm requiring an undirected graph on a directed graph, the program reports that an undirected graph must be generated; see Fig.~\ref{fig:res-error-2}.

\begin{figure}[H]
	\centering
	\includegraphics[width=0.9\textwidth,height=0.45\textheight,keepaspectratio]{\detokenize{error handling/error handle2 for 16 eulerian check.png}}
	\caption{Handling an incorrect graph type for the Eulerian algorithm}
	\label{fig:res-error-2}
\end{figure}

The correctness of vertex numbers and the admissibility of source and sink selection are also checked. With incorrect input, the program does not terminate abnormally; it prints a clear message and returns the user to the menu; see Fig.~\ref{fig:res-error-3}.

\begin{figure}[H]
	\centering
	\includegraphics[width=0.9\textwidth,height=0.45\textheight,keepaspectratio]{\detokenize{error handling/error handle3 .png}}
	\caption{Handling incorrect user input}
	\label{fig:res-error-3}
\end{figure}

\clearpage
\section*{Conclusion}
\addcontentsline{toc}{section}{Conclusion}

Based on the completed laboratory works, the following conclusions can be drawn:

In the first laboratory work, an algorithm for generating a random connected acyclic graph using the normal distribution was implemented. For the obtained graph, vertex eccentricities, radius, diameter, center, and diametral vertices were calculated. Shimbell's method for finding minimum and maximum paths of a specified length was also implemented, as well as an algorithm for determining all routes between two selected vertices.

In the second laboratory work, depth-first graph traversal was performed. For a weighted graph, the Floyd-Warshall algorithm was implemented; it finds shortest paths between all pairs of vertices and reconstructs the path between vertices selected by the user. The program also outputs the number of iterations, which makes it possible to estimate the complexity of the executed algorithms.

In the third laboratory work, capacity and cost matrices were generated based on a directed graph. For the constructed transport network, the maximum flow was found using the Ford-Fulkerson algorithm. After that, the minimum-cost flow was computed for a value equal to $\lfloor 2/3 \cdot \varphi_{\max} \rfloor$, using the previously implemented Floyd-Warshall algorithm to find shortest paths by cost.

In the fourth laboratory work, the number of spanning trees was found for an undirected graph using Kirchhoff's matrix-tree theorem. The minimum spanning tree was built using Boruvka's algorithm. The obtained spanning tree was encoded by the Prufer code and then decoded while preserving edge weights. The search for a minimum edge cover was also implemented both for the original graph and for the constructed spanning tree.

In the fifth laboratory work, the graph was checked for Eulerian property. If the graph was not Eulerian, the program modified it with output of the performed changes and then built an Eulerian cycle. For the even-numbered variant, the fundamental cutset system based on the minimum spanning tree was implemented, as well as obtaining other cutsets using the symmetric difference of selected fundamental cutsets.

The completion of the laboratory works made it possible to study and implement the main graph theory algorithms for directed and undirected graphs: traversals, shortest paths, flows, spanning trees, covers, tree coding, Eulerian cycles, and fundamental cutsets. All algorithms are combined in one console program with a menu, which allows the user to sequentially generate a graph, build the required matrices, and apply algorithms from different laboratory works.

An advantage of the program is its modular structure: common graph entities, distribution generation, and weight matrix generation are placed in separate files, while algorithms are divided by laboratory works. The use of STL containers, especially \texttt{vector}, simplifies work with matrices and edge lists. The program includes checks for user input and restrictions on launching algorithms, for example the requirement of a directed graph for flow algorithms and an undirected graph for spanning trees, Eulerian cycles, and cutsets.

The disadvantages include the absence of a graphical graph representation: results are printed to the console as matrices, edge lists, and textual algorithm logs. In addition, some algorithms do not have the most optimal asymptotic complexity for large graphs, for example the minimum-cost flow search through repeated application of the Floyd-Warshall algorithm. In the future, the program can be extended with graph visualization and more efficient implementations for large graphs.

The laboratory works were implemented in C++ in the CLion development environment.

\clearpage
\addcontentsline{toc}{section}{References}
\begin{thebibliography}{99}

	\bibitem{novikov}
	Novikov F.\,A. Discrete Mathematics for Programmers~/ F.\,A.~Novikov.~---
	3rd ed.~--- St. Petersburg:~Piter, 2009.~--- 364~p.

	\bibitem{shaporev}
	Shaporev S.\,D. Discrete Mathematics~/ S.\,D.~Shaporev.~---
	St. Petersburg:~BHV-Petersburg, 2010.~--- 400~p.

	\bibitem{romanovsky}
	Romanovsky I.\,V. Discrete Analysis~/ I.\,V.~Romanovsky.~---
	3rd ed.~--- St. Petersburg:~Nevsky Dialect; BHV-Petersburg, 2003.~--- 320~p.

	\bibitem{haggarty}
	Haggarty R. Discrete Mathematics for Programmers~/ R.~Haggarty; translated from English.~---
	Moscow:~Technosphere, 2005.~--- 326~p.

	\bibitem{cormen}
	Cormen T., Leiserson C., Rivest R., Stein C. Introduction to Algorithms~/ translated from English.~---
	3rd ed.~--- Moscow:~Williams, 2013.~--- 1328~p.

	\bibitem{asanov}
	Asanov M.\,O., Baransky V.\,A., Rasin V.\,V. Discrete Mathematics: Graphs, Matroids, Algorithms~/ M.\,O.~Asanov, V.\,A.~Baransky, V.\,V.~Rasin.~---
	2nd ed.~--- St. Petersburg:~Lan, 2010.~--- 368~p.

	\bibitem{vadzinsky}
	Vadzinsky R.\,N. Handbook of Probability Distributions~/ R.\,N.~Vadzinsky.~---
	St. Petersburg:~Nauka, 2001.~--- 295~p.

	\bibitem{erickson}
	Erickson J. Algorithms~/ J.~Erickson; translated from English.~---
	Moscow:~Williams, 2022.~--- 960~p.

	\bibitem{telematika}
	Section "Telematics"~/ electronic text.~---
	URL: \url{https://tema.spbstu.ru/tgraph/}
	(accessed: 20.05.2026)

\end{thebibliography}

\clearpage
\section*{Appendix A. Common Modules and the Combined Program}
\addcontentsline{toc}{section}{Appendix A. Common Modules and the Combined Program}

This appendix contains the complete source code of the common modules used by all laboratory works, as well as the code of the combined console program with a menu.

\lstinputlisting[language=C++, caption={File constants.h}, label={lst:shared-constants-h}]{../graph-theory-labs/shared/constants.h}

\lstinputlisting[language=C++, caption={File graph.h}, label={lst:shared-graph-h}]{../graph-theory-labs/shared/graph.h}

\lstinputlisting[language=C++, caption={File graph.cpp}, label={lst:shared-graph-cpp}]{../graph-theory-labs/shared/graph.cpp}

\lstinputlisting[language=C++, caption={File distribution.h}, label={lst:shared-distribution-h}]{../graph-theory-labs/shared/distribution.h}

\lstinputlisting[language=C++, caption={File distribution.cpp}, label={lst:shared-distribution-cpp}]{../graph-theory-labs/shared/distribution.cpp}

\lstinputlisting[language=C++, caption={File dag\_generator.h}, label={lst:shared-dag-generator-h}]{../graph-theory-labs/shared/dag_generator.h}

\lstinputlisting[language=C++, caption={File dag\_generator.cpp}, label={lst:shared-dag-generator-cpp}]{../graph-theory-labs/shared/dag_generator.cpp}

\lstinputlisting[language=C++, caption={File weight\_matrix.h}, label={lst:shared-weight-matrix-h}]{../graph-theory-labs/shared/weight_matrix.h}

\lstinputlisting[language=C++, caption={File weight\_matrix.cpp}, label={lst:shared-weight-matrix-cpp}]{../graph-theory-labs/shared/weight_matrix.cpp}

\lstinputlisting[language=C++, caption={File main.cpp of the combined program}, label={lst:lab-combined-main-cpp}]{../graph-theory-labs/lab_combined/main.cpp}

\clearpage
\section*{Appendix B. Source Code of Laboratory Work No. 1}
\addcontentsline{toc}{section}{Appendix B. Source Code of Laboratory Work No. 1}

\lstinputlisting[language=C++, caption={File eccentricity.h}, label={lst:lab1-eccentricity-h}]{../graph-theory-labs/lab1/eccentricity.h}

\lstinputlisting[language=C++, caption={File eccentricity.cpp}, label={lst:lab1-eccentricity-cpp}]{../graph-theory-labs/lab1/eccentricity.cpp}

\lstinputlisting[language=C++, caption={File shimbell.h}, label={lst:lab1-shimbell-h}]{../graph-theory-labs/lab1/shimbell.h}

\lstinputlisting[language=C++, caption={File shimbell.cpp}, label={lst:lab1-shimbell-cpp}]{../graph-theory-labs/lab1/shimbell.cpp}

\lstinputlisting[language=C++, caption={File path\_counter.h}, label={lst:lab1-path-counter-h}]{../graph-theory-labs/lab1/path_counter.h}

\lstinputlisting[language=C++, caption={File path\_counter.cpp}, label={lst:lab1-path-counter-cpp}]{../graph-theory-labs/lab1/path_counter.cpp}

\clearpage
\section*{Appendix C. Source Code of Laboratory Work No. 2}
\addcontentsline{toc}{section}{Appendix C. Source Code of Laboratory Work No. 2}

\lstinputlisting[language=C++, caption={File dfs.h}, label={lst:lab2-dfs-h}]{../graph-theory-labs/lab2/dfs.h}

\lstinputlisting[language=C++, caption={File dfs.cpp}, label={lst:lab2-dfs-cpp}]{../graph-theory-labs/lab2/dfs.cpp}

\lstinputlisting[language=C++, caption={File floyd\_warshall.h}, label={lst:lab2-floyd-warshall-h}]{../graph-theory-labs/lab2/floyd_warshall.h}

\lstinputlisting[language=C++, caption={File floyd\_warshall.cpp}, label={lst:lab2-floyd-warshall-cpp}]{../graph-theory-labs/lab2/floyd_warshall.cpp}

\clearpage
\section*{Appendix D. Source Code of Laboratory Work No. 3}
\addcontentsline{toc}{section}{Appendix D. Source Code of Laboratory Work No. 3}

\lstinputlisting[language=C++, caption={File ford\_fulkerson.h}, label={lst:lab3-ford-fulkerson-h}]{../graph-theory-labs/lab3/ford_fulkerson.h}

\lstinputlisting[language=C++, caption={File ford\_fulkerson.cpp}, label={lst:lab3-ford-fulkerson-cpp}]{../graph-theory-labs/lab3/ford_fulkerson.cpp}

\lstinputlisting[language=C++, caption={File min\_cost\_flow.h}, label={lst:lab3-min-cost-flow-h}]{../graph-theory-labs/lab3/min_cost_flow.h}

\lstinputlisting[language=C++, caption={File min\_cost\_flow.cpp}, label={lst:lab3-min-cost-flow-cpp}]{../graph-theory-labs/lab3/min_cost_flow.cpp}

\clearpage
\section*{Appendix E. Source Code of Laboratory Work No. 4}
\addcontentsline{toc}{section}{Appendix E. Source Code of Laboratory Work No. 4}

\lstinputlisting[language=C++, caption={File kirchhoff.h}, label={lst:lab4-kirchhoff-h}]{../graph-theory-labs/lab4/kirchhoff.h}

\lstinputlisting[language=C++, caption={File kirchhoff.cpp}, label={lst:lab4-kirchhoff-cpp}]{../graph-theory-labs/lab4/kirchhoff.cpp}

\lstinputlisting[language=C++, caption={File boruvka.h}, label={lst:lab4-boruvka-h}]{../graph-theory-labs/lab4/boruvka.h}

\lstinputlisting[language=C++, caption={File boruvka.cpp}, label={lst:lab4-boruvka-cpp}]{../graph-theory-labs/lab4/boruvka.cpp}

\lstinputlisting[language=C++, caption={File edge\_cover.h}, label={lst:lab4-edge-cover-h}]{../graph-theory-labs/lab4/edge_cover.h}

\lstinputlisting[language=C++, caption={File edge\_cover.cpp}, label={lst:lab4-edge-cover-cpp}]{../graph-theory-labs/lab4/edge_cover.cpp}

\lstinputlisting[language=C++, caption={File prufer.h}, label={lst:lab4-prufer-h}]{../graph-theory-labs/lab4/prufer.h}

\lstinputlisting[language=C++, caption={File prufer.cpp}, label={lst:lab4-prufer-cpp}]{../graph-theory-labs/lab4/prufer.cpp}

\clearpage
\section*{Appendix F. Source Code of Laboratory Work No. 5}
\addcontentsline{toc}{section}{Appendix F. Source Code of Laboratory Work No. 5}

\lstinputlisting[language=C++, caption={File eulerian.h}, label={lst:lab5-eulerian-h}]{../graph-theory-labs/lab5/eulerian.h}

\lstinputlisting[language=C++, caption={File eulerian.cpp}, label={lst:lab5-eulerian-cpp}]{../graph-theory-labs/lab5/eulerian.cpp}

\lstinputlisting[language=C++, caption={File fundamental\_cutsets.h}, label={lst:lab5-fundamental-cutsets-h}]{../graph-theory-labs/lab5/fundamental_cutsets.h}

\lstinputlisting[language=C++, caption={File fundamental\_cutsets.cpp}, label={lst:lab5-fundamental-cutsets-cpp}]{../graph-theory-labs/lab5/fundamental_cutsets.cpp}

\end{document}
